\section{실패 유형 판정 규칙}
\begin{table}[H]
\centering\small
\caption{12종 실패 유형 체계}
\begin{tabularx}{\linewidth}{@{}p{24mm}p{32mm}Y@{}}
\toprule
분류 & 유형 & 판정 규칙 (요약) \\
\midrule
\multirow{6}{24mm}{초기 환경·물리 6종} & 충돌 & 로봇 몸체와 장애물의 접촉이 지속되며 작업 실패 \\
& 도달 불가 & 목표점 계획이 반복 거절되거나, 물체 운반 후 목표 위치 미완료 \\
& 장애물 간섭 & 마지막 10초에 경로 계획 거절 2회 이상, 장애물과 0.55\,m 이내, 정체 \\
& 목표 위치 점유 & 목표 영역이 물체로 점유된 상태에서 종료 구간 정체 \\
& 작업자 안전 위험 & 사람과 접촉 또는 1\,m 이내 0.5초 이상 연속 노출 \\
& 인식 오류 & 잘못된 대상 선택, 가림에 의한 인식 실패 \\
\midrule
\multirow{4}{24mm}{행동 수행 4종} & 넘어짐 & 몸통 높이 0.45\,m 미만 또는 기울기 60도 이상이 0.5초 이상 \\
& 이동 정체 & 2초 이상의 이동 실행에서 순이동 0.05\,m 이하 \\
& 물체 조작 실패 & 상자 밀기 등 조작 요청이 최종적으로 거절·실패 \\
& 목표 완료 실패 & 목표 반경 안에 도착했으나 체류 시간 미충족 \\
\midrule
\multirow{2}{24mm}{시스템 관측 2종} & 최종 목표점 계획 거절 & 경로 계획기가 최종 목표점을 막힘으로 2회 이상 거절, 마지막 행동도 거절 \\
& 중간 목표 반복 & 최종 목표 밖 같은 중간 목표에 반복 도착 \\
\bottomrule
\end{tabularx}
\end{table}
한 실행에서 여러 현상이 관측되면 대표 유형 1개만 집계하고 나머지는 보조 진단으로 보존한다.
상태 기록 간격이 0.25초를 넘거나 필수 기록이 없으면 판정을 보류한다.
